\documentclass[12pt]{article}

%----------------- Packages ------------------
\usepackage[utf8]{inputenc}
\usepackage[T1]{fontenc}
\usepackage{amsmath, amssymb}
\usepackage{geometry}
\usepackage{xcolor}
\usepackage{titlesec}
\usepackage{fancyhdr}
\usepackage[breakable]{tcolorbox}
\usepackage{graphicx}
\usepackage{tikz}
\usepackage{float}
\usetikzlibrary{arrows.meta, decorations.markings}

%----------------- Page Setup -----------------
\geometry{a4paper, margin=2.5cm}
\pagestyle{fancy}
\fancyhf{}
\rhead{Final Exam — \textbf{2019}}
\lhead{Probability and Statistics}
\cfoot{\thepage}

%----------------- Title Styling --------------
\titleformat{\section}{\normalfont\Large\bfseries}{Exercise \thesection:}{1em}{}
\titleformat{\subsection}{\normalfont\bfseries}{Answer:}{1em}{}

%----------------- Custom Boxes ----------------
\tcbuselibrary{listingsutf8}
\newtcolorbox{answerbox}{
  colback=gray!10,
  colframe=black,
  fonttitle=\bfseries,
  title=Answer Area,
  breakable,
  before skip=10pt,
  after skip=10pt
}

%----------------- Document Start --------------
\begin{document}

%----------------- Exam Info -------------------
\begin{center}
  \Large\textbf{UNIVERSITY NAME} \\[1em]
  \large\textit{Probability and Statistics — Final Exam} \\[0.5em]
  \large\textit{Year: 2019} \\[2em]
\end{center}

\vspace{0.5cm}

%----------------- EXERCISE 1 ------------------
\section{}
Let \( X \) be a random variable uniformly distributed on \([0,1]\) with probability density function  
\[
f_X(x) = 
\begin{cases} 
1 & \text{if } 0 \leq x \leq 1, \\
0 & \text{otherwise.}
\end{cases}
\]  
Define \( Y = -\ln X \).

\begin{enumerate}
    \item Find the probability density function of \( Y \).
    \item Compute \( P\left( |Y - 1| > \frac{1}{10} \right) \).
\end{enumerate}

\begin{answerbox}
\begin{enumerate}
  \item Probability density function of \( Y \)

Since \( X \sim \mathcal{U}([0,1]) \), and \( Y = -\ln X \). The transformation is strictly decreasing.
The CDF method:
For \( y > 0 \),
\[
F_Y(y) = P(Y \le y) = P(-\ln X \le y) = P(X \ge e^{-y}) = 1 - F_X(e^{-y}).
\]
Since \( f_X(x)=1 \) for \( x \in [0,1] \), \( F_X(x) = x \) for \( 0 \le x \le 1 \). For \( y>0 \), \( e^{-y} \in (0,1] \), so:
\[
F_Y(y) = 1 - e^{-y}.
\]
Differentiating:
\[
f_Y(y) = \frac{d}{dy}(1 - e^{-y}) = e^{-y}, \quad y > 0.
\]
Thus,
\[
\boxed{f_Y(y) = e^{-y} \mathbf{1}_{(0,\infty)}(y)}.
\]

\item Compute \( P\left( |Y - 1| > \frac{1}{10} \right) \)
\[
P\left( |Y - 1| > \frac{1}{10} \right) = 1 - P\left( |Y - 1| \le \frac{1}{10} \right) = 1 - P\left( 0.9 \le Y \le 1.1 \right).
\]
\[
P(0.9 \le Y \le 1.1) = \int_{0.9}^{1.1} e^{-y} \, dy = \left[ -e^{-y} \right]_{0.9}^{1.1} = e^{-0.9} - e^{-1.1}.
\]
Thus,
\[
P\left( |Y - 1| > \frac{1}{10} \right) = 1 - \left( e^{-0.9} - e^{-1.1} \right) = 1 - e^{-0.9} + e^{-1.1}.
\]
\[
\boxed{1 - e^{-0.9} + e^{-1.1}}.
\]
\end{enumerate}
\end{answerbox}

%----------------- EXERCISE 2 ------------------
\section{}
A player starts a video game and plays several matches. It is assumed that the probability of winning the first match is \( 0.1 \):  

- If they win a match, the probability of winning the next one is \( 0.8 \).  

- If they lose a match, the probability of winning the next one is \( 0.6 \).  

Let \( G_n \) be the event “the player wins the \( n \)-th match” and let \( p_n = P(G_n) \).

\begin{enumerate}
    \item Calculate \( p_2 \).
    \item Given that the player won the second match, calculate the probability that they lost the first one.
    \item Find the probability that the player wins at least one match in the first three matches.
    \item Show that for all natural numbers \( n \geq 1 \), we have  
    \[
    p_{n+1} = \frac{1}{5} p_n + \frac{3}{5}.
    \]
    \item Assume the sequence \( (p_n)_{n \geq 1} \) converges to a limit \( \ell \). Determine \( \ell \) and show that the sequence \( (p_n - \ell)_{n \geq 1} \) is geometric.
    \item Deduce an expression for \( p_n \) as a function of \( n \).
\end{enumerate}

\begin{answerbox}
Given: \( P(G_1) = p_1 = 0.1 \), and
\[
P(G_{n+1} \mid G_n) = 0.8, \quad P(G_{n+1} \mid G_n^c) = 0.6.
\]

\begin{enumerate}
  \item Calculate \( p_2 \)
\[
p_2 = P(G_2) = P(G_2 \mid G_1)P(G_1) + P(G_2 \mid G_1^c)P(G_1^c)
= 0.8 \cdot 0.1 + 0.6 \cdot 0.9  = 0.62.
\]

  \item Given \( G_2 \), probability they lost the first one
\[
P(G_1^c \mid G_2) = \frac{P(G_2 \mid G_1^c) P(G_1^c)}{P(G_2)} = \frac{0.6 \cdot 0.9}{0.62} = \frac{0.54}{0.62} = \frac{27}{31}.
\]

  \item Probability that the player wins at least one match in the first three
Easier to compute complement: loses all three matches.
Let \( L_i = G_i^c \). Then:
$$
P(L_1) = 0.9, \quad P(L_2 \mid L_1) = 1 - 0.6 = 0.4, \quad \\
$$
$$
P(L_3 \mid L_2) = 1 - 0.6 = 0.4 \quad (\text{since after loss, win prob } 0.6)
$$
But careful: the loss-loss transition probability is actually \( P(L_{n+1} \mid L_n) = 1 - 0.6 = 0.4 \), yes.

So:
\[
P(L_1 \cap L_2 \cap L_3) = P(L_1) P(L_2 \mid L_1) P(L_3 \mid L_2) = 0.9 \cdot 0.4 \cdot 0.4 = 0.144.
\]
Thus:
\[
P(\text{at least one win}) = 1 - 0.144 = 0.856.
\]

  \item Show that \( p_{n+1} = \frac{1}{5} p_n + \frac{3}{5} \)
\[
p_{n+1} = P(G_{n+1}) = P(G_{n+1} \mid G_n) P(G_n) + P(G_{n+1} \mid G_n^c) P(G_n^c)
\]
\[
= 0.8 p_n + 0.6 (1 - p_n) = 0.8 p_n + 0.6 - 0.6 p_n = 0.2 p_n + 0.6 = \frac{1}{5} p_n + \frac{3}{5}.
\]

  \item Limit \( \ell \) and show \( (p_n - \ell) \) geometric
If limit \( \ell \) exists, then from recurrence: \( \ell = \frac{1}{5} \ell + \frac{3}{5} \) \(\Rightarrow\) \( \ell - \frac{1}{5}\ell = \frac{3}{5} \) \(\Rightarrow\) \( \frac{4}{5}\ell = \frac{3}{5} \) \(\Rightarrow\) \( \ell = \frac{3}{4} \).
Now consider \( d_n = p_n - \ell = p_n - \frac{3}{4} \). Then:
\[
d_{n+1} = p_{n+1} - \frac{3}{4} = \left( \frac{1}{5} p_n + \frac{3}{5} \right) - \frac{3}{4} = \frac{1}{5} p_n - \frac{3}{20}.
\]
But note \( \frac{3}{20} = \frac{1}{5} \cdot \frac{3}{4} \), so:
\[
d_{n+1} = \frac{1}{5} p_n - \frac{1}{5} \cdot \frac{3}{4} = \frac{1}{5} \left( p_n - \frac{3}{4} \right) = \frac{1}{5} d_n.
\]
Thus \( (d_n) \) is geometric with ratio \( \frac{1}{5} \).

\[
\ell = \frac{3}{4}, \quad p_n - \ell \text{ is geometric with ratio } \frac{1}{5}.
\]

  \item Expression for \( p_n \)
Since \( d_1 = p_1 - \frac{3}{4} = 0.1 - 0.75 = -0.65 = -\frac{13}{20} \), then:
\[
d_n = d_1 \left( \frac{1}{5} \right)^{n-1} = -\frac{13}{20} \cdot \left( \frac{1}{5} \right)^{n-1}.
\]
Thus:
\[
p_n = \ell + d_n = \frac{3}{4} - \frac{13}{20} \cdot \left( \frac{1}{5} \right)^{n-1}.
\]
\end{enumerate}
\end{answerbox}


%----------------- END ------------------
\end{document}